\documentclass[letterpaper,11pt,leqno]{article}
\usepackage{paper}
\usepackage{natbib}
\bibliographystyle{bibliography}

\usepackage{xcolor}
\usepackage{hyperref}
\usepackage{threeparttable}
\usepackage{float}
\usepackage{booktabs}

\definecolor{customblue}{rgb}{0,0.337,0.702}

\hypersetup{
    colorlinks   = true,
    citecolor    = customblue,
    linkcolor    = customblue,
    urlcolor     = customblue
}

\newenvironment{fignotes2}{\begin{quote}\setstretch{1.2}\scriptsize}{\end{quote}}

% Every number quoted in the text, generated by code/6. note_numbers.r
% Generated by code/6. note_numbers.r -- do not edit by hand.
% ADRH 2023 cross-section, nowcast to 2025; tract GB2 fits.
\newcommand{\BaseYear}{2023}
\newcommand{\TargetYear}{2025}
\newcommand{\TractsTotal}{37,072}
\newcommand{\TractPopMean}{1,315}
\newcommand{\TractPopSD}{678}
\newcommand{\TractEquivMean}{22,036}
\newcommand{\TractPcMean}{15,025}
\newcommand{\EquivOverPcPct}{47}
\newcommand{\TractGiniMean}{28.6}
\newcommand{\TractDepRatioMean}{0.63}
\newcommand{\TractMeanAge}{45.6}
\newcommand{\TractSingleHhMean}{30.7}
\newcommand{\MissEquivPct}{5.5}
\newcommand{\GiniMissTracts}{2,047}
\newcommand{\GiniMissPct}{5.5}
\newcommand{\GiniMissPopPct}{0.3}
\newcommand{\GiniMissMedianPop}{57}
\newcommand{\ObsGiniMinPop}{100}
\newcommand{\GiniImputedTracts}{1,370}
\newcommand{\TractsDropped}{677}
\newcommand{\LocMeanTracts}{1,370}
\newcommand{\MissEquivAfterPct}{1.8}
\newcommand{\ShareTracts}{31,629}
\newcommand{\ShareTractsPopPct}{96}
\newcommand{\MedBinnedPct}{98.5}
\newcommand{\CapMedian}{44,450}
\newcommand{\CapTracts}{148}
\newcommand{\CapPopPct}{0.5}
\newcommand{\CapGini}{44.2}
\newcommand{\CapGiniTracts}{182}
\newcommand{\AdrhNatMedian}{19,250}
\newcommand{\AdrhNatMean}{22,318}
\newcommand{\AdrhNatGini}{32.1}
\newcommand{\AdrhNatPEightyTwenty}{2.8}
\newcommand{\MunAdrh}{8,139}
\newcommand{\MunLookup}{8,058}
\newcommand{\MunDropped}{81}
\newcommand{\MunSingleTract}{5,805}
\newcommand{\MunSingleTractPct}{72}
\newcommand{\MunSingleTractPopPct}{5.9}
\newcommand{\MunSynthetic}{1,326}
\newcommand{\MunSyntheticPct}{16}
\newcommand{\MunSyntheticPopPct}{0.28}
\newcommand{\MunSyntheticIncomeValues}{84}
\newcommand{\GbTracts}{35,025}
\newcommand{\GbShareTracts}{31,629}
\newcommand{\GbPriorTracts}{3,396}
\newcommand{\GbPriorTrain}{30,665}
\newcommand{\GbGapMuns}{47}
\newcommand{\GbPriorRMSE}{1.1}
\newcommand{\GbConvergedPct}{99.98}
\newcommand{\TolShare}{0.4}
\newcommand{\TolGini}{0.25}
\newcommand{\TolMeanPct}{1}
\newcommand{\TolPPct}{1}
\newcommand{\MedPullPct}{2}
\newcommand{\MedOutPct}{0.3}
\newcommand{\CapAllowPct}{10}
\newcommand{\ValTracts}{35,025}
\newcommand{\FitMedInBinPct}{83}
\newcommand{\FitMeanAbsErr}{1.0}
\newcommand{\FitGiniAbsErr}{0.19}
\newcommand{\FitPAbsErr}{2.6}
\newcommand{\ValSkewSlope}{0.96}
\newcommand{\ValSkewRsq}{0.88}
\newcommand{\ValPSlope}{0.96}
\newcommand{\ValPRsq}{0.94}
\newcommand{\FitShareMAE}{0.51}
\newcommand{\HoldShareMAE}{0.71}
\newcommand{\HoldShareMAEMin}{0.55}
\newcommand{\HoldShareMAEMax}{0.89}
\newcommand{\RootErrExp}{-9}
\newcommand{\AppNatMedian}{20,975}
\newcommand{\MixNatMedian}{19,029}
\newcommand{\MixNatMedianErr}{$-$1.1}
\newcommand{\AdrhNatMedianLo}{18,900}
\newcommand{\AdrhNatMedianHi}{19,600}
\newcommand{\MixNatMean}{22,326}
\newcommand{\MixNatMeanErr}{$+$0.04}
\newcommand{\MixNatGini}{32.2}
\newcommand{\MixNatPEightyTwenty}{2.71}
\newcommand{\MixNatPEightyTwentyBinned}{2.8}
\newcommand{\MixPOne}{2,492}
\newcommand{\ProvN}{52}
\newcommand{\ProvMedAbsErr}{0.8}
\newcommand{\ProvMedInBinPct}{95}
\newcommand{\ProvGiniAbsErr}{0.07}
\newcommand{\ProvPAbsErr}{1.5}
\newcommand{\ProvPBias}{$-$0.3}
\newcommand{\MunValN}{2,253}
\newcommand{\MunMedAbsErr}{0.9}
\newcommand{\MunMedInBinPct}{89}
\newcommand{\MunGiniAbsErr}{0.11}
\newcommand{\MunPAbsErr}{2.0}
\newcommand{\MunPBias}{$+$0.5}
\newcommand{\ThrFiveKObs}{4.3}
\newcommand{\ThrFiveKModel}{3.9}
\newcommand{\ThrOtherMaxDiff}{0.3}
\newcommand{\ThrAllMaxDiff}{0.4}
\newcommand{\ThrPovEuro}{11,550}
\newcommand{\ThrPovObs}{21.3}
\newcommand{\ThrPovModel}{21.3}
\newcommand{\ThrPovModelP}{21}
\newcommand{\ThrPovObsP}{21}
\newcommand{\ThrTopEuro}{38,500}
\newcommand{\ThrTopObs}{10.0}
\newcommand{\ThrTopModel}{10.2}
\newcommand{\CvFolds}{5}
\newcommand{\XgbRounds}{100}
\newcommand{\XgbDepth}{6}
\newcommand{\XgbEta}{0.3}
\newcommand{\GiniTrainTracts}{35,025}
\newcommand{\GiniSD}{3.87}
\newcommand{\GiniObsMean}{28.6}
\newcommand{\GiniPredMean}{30.2}
\newcommand{\OlsRMSE}{2.94}
\newcommand{\XgbRMSE}{2.57}
\newcommand{\XgbRsq}{0.56}
\newcommand{\XgbRMSEGain}{12.8}
\newcommand{\SmallTrainMin}{300}
\newcommand{\SmallTestMaxPop}{299}
\newcommand{\SmallTestTracts}{1,807}
\newcommand{\XgbSmallRMSE}{4.30}
\newcommand{\NowcastFirstYear}{2016}
\newcommand{\NowcastLastYear}{2023}
\newcommand{\NowcastRho}{0.93}
\newcommand{\NowcastMidYear}{2024}
\newcommand{\NowcastGrowthMid}{5.8}
\newcommand{\NowcastGrowthTarget}{4.2}
\newcommand{\NowcastGrowth}{10.2}
\newcommand{\AeatNetTarget}{4.5}
\newcommand{\AeatGrossTarget}{5.3}
\newcommand{\RhoLastThree}{0.92}
\newcommand{\RhoExclCovid}{0.91}
\newcommand{\RhoOrigin}{0.91}
\newcommand{\NowcastGrowthMin}{10.0}
\newcommand{\NowcastGrowthMax}{11.0}
\newcommand{\NowcastBtN}{5}
\newcommand{\NowcastBtFirst}{2017}
\newcommand{\NowcastBtLast}{2021}
\newcommand{\NowcastBtMAE}{0.5}
\newcommand{\NowcastBtMax}{1.0}
\newcommand{\NowcastBtOneMAE}{0.3}
\newcommand{\VdWithin}{84.6}
\newcommand{\VdBetweenTract}{8.8}
\newcommand{\VdBetweenMun}{2.7}
\newcommand{\VdBetweenProv}{0.5}
\newcommand{\VdBetweenCcaa}{3.5}
\newcommand{\VdCcaaN}{17}
\newcommand{\VdCcaaWithin}{90}
\newcommand{\VdCcaaTract}{7}
\newcommand{\VdCcaaRest}{3}
\newcommand{\VdCcaaWithinMin}{80}
\newcommand{\VdCcaaWithinMax}{94}


\begin{document}

\title{\Large \textsc{Estimating Spain's income distribution from neighborhood-level data}}

\author{\href{https://pablogguz.github.io/}{\large{Pablo García-Guzmán}\thanks{Email: \href{mailto:garciagp@ebrd.com}{garciagp@ebrd.com}. I thank Pedro Torres for his valuable feedback and comments. Agentic AI tools were used to write and debug the code of the data pipeline and the web application, to review the methodology and to edit the text. All remaining errors are my own. The views expressed in this methodology note are my own and do not reflect the views of the EBRD. The code and replication package for this note are available at \href{https://github.com/pablogguz/compara_tu_ingreso_app}{github.com/pablogguz/compara\_tu\_ingreso\_app}.}} \\ \normalsize{}}

\date{\normalsize \today \\}

\maketitle
\begin{abstract}
\onehalfspacing
This methodology note develops a transparent and simple framework to estimate aggregate income distributions in Spain exploiting the availability of granular inequality metrics. Specifically, I leverage distributional indicators derived from personal income-tax records at the census tract-level to reconstruct local income distributions assuming a generalized beta distribution of the second kind (GB2). National, regional and municipal income distributions can be subsequently derived as population-weighted mixtures of these tract-level distributions. This approach leverages well-documented regularities in how incomes are distributed within demographically homogeneous groups, and the resulting estimates closely match observed distributional statistics in validation tests: the national shares of the population below nine published income thresholds are reproduced within \ThrAllMaxDiff{} percentage points. While the focus is on Spain, the approach uses only published small-area statistics and can be readily adapted to other countries with comparable data available. The framework is implemented in \href{https://comparatuingreso.es/}{\mbox{comparatuingreso.es}}, a publicly available web platform that enables Spanish households to calculate their relative position within the income distribution.
\end{abstract}

\thispagestyle{empty}

\newpage

\setcounter{page}{1}

\section{Introduction}

Understanding the income distribution is crucial for economic analysis and policy design, yet its measurement remains challenging. Statistical offices increasingly publish income statistics for small geographic units, such as mean and median incomes, inequality indices and the shares of the population below given income levels, but rarely the full distributions behind them, and how to recover those distributions from published summary statistics has received little attention. In this note, I propose a transparent and simple method to do so. The approach builds on two insights: first, that the income distribution of a small, demographically homogeneous area can be well described by a flexible parametric family fitted to the statistics published for it; and second, that the distribution of any larger area can then be obtained as a population-weighted mixture of these local distributions.

I apply the method to Spain, where the Spanish Statistical Office publishes a set of distributional indicators derived from personal income-tax records for each of its \TractsTotal{} census tracts, the smallest areas for which official statistics are released, with \TractPopMean{} residents on average. I model each tract's distribution as a generalized beta distribution of the second kind (GB2), a four-parameter family that nests the log-normal as a limiting case, fitted to these indicators. The resulting national, provincial and municipal distributions closely match the statistics published for those levels, which are not used in the estimation: the national median falls inside its published interval, the Gini coefficients of provinces and municipalities are reproduced within \ProvGiniAbsErr{} and \MunGiniAbsErr{} points on average, and the national shares of the population below nine published income thresholds within \ThrAllMaxDiff{} percentage points. The method is implemented in \href{https://comparatuingreso.es/}{\mbox{comparatuingreso.es}}, a publicly available web platform that enables Spanish households to calculate their relative position within the income distribution of Spain, their province and their municipality.

The rest of this note is organized as follows. Section \ref{sec:data} describes the data sources, their coverage and scope. Section \ref{sec:method} outlines the methodology, including the modeling of local income distributions and their aggregation. Section \ref{sec:results} presents the validation results and a decomposition of income inequality. Finally, Section \ref{sec:conc} concludes.

\section{Data}\label{sec:data}


\paragraph{Income} The main source is the INE's Household Income Distribution Atlas (\textit{Atlas de Distribución de Renta de los Hogares}, ADRH; \citealt{ine2025adrh}). The ADRH links personal income-tax records from the national tax agency (AEAT) and the foral tax authorities of the Basque Country and Navarre to the population register (\textit{Fichero Precensal de Población}) and attributes income to the census tract where each recipient lives. Income refers to a calendar year and population to 1 January of the following year; people living in collective establishments, such as care homes or barracks, are excluded. I use the latest edition, which refers to income in \BaseYear{}.

Net income is gross income minus the personal income tax due, employees' social security contributions and the wealth tax. Gross income comprises wages, including remuneration in kind and employers' contributions to pension plans; pensions; unemployment benefits; other public benefits, including tax-exempt ones such as disability, dependency and family benefits and regional minimum-income schemes; interest, dividends and other capital income; rental income; business income; and the notional income that the income tax attributes to second homes. Non-filers are covered through the information returns of those who pay them. The ADRH excludes income that the tax authorities do not record, such as informal earnings, transfers between households and foreign pensions, as well as capital gains. It also covers only people living in households with some recorded income, which in 2016 included more than 98.6\% of the population of the common fiscal territory.

Household net income is divided by the household's consumption units on the modified OECD scale (1 for the first adult, 0.5 for every other member aged 14 or over and 0.3 for every member under 14) and the result is assigned to every member. The ADRH publishes thirteen statistics of this distribution of persons: its mean, median, Gini coefficient and P80/P20 ratio, and the shares of the population with income below three fixed amounts (\texteuro 5,000, \texteuro 7,500 and \texteuro 10,000) and below or above six fractions of the national median (40, 50, 60, 140, 160 and 200\%). All the distributions in this note are distributions of persons.

\paragraph{Publication rules} Four features of the ADRH matter for what follows. First, the mean and median of income per consumption unit, the Gini coefficient and the P80/P20 ratio are published only for tracts and municipalities with at least 100 residents, and the shares only for those with at least 500 (\ShareTracts{} tracts, with \ShareTractsPopPct{}\% of the population). Below that size only income per person and per household is published, taken from the district or municipality that contains the tract or, for municipalities with fewer than 100 residents, averaged over the municipalities of that size in the same agricultural district (\textit{comarca agraria}). Second, the Gini coefficient is computed after replacing the lowest and highest 1.5\% of incomes with the nearest remaining value. Third, published tract indicators are bounded at the population-weighted 0.1st and 99.5th percentiles across tracts; in \BaseYear{} this caps the median of \CapTracts{} tracts, home to \CapPopPct{}\% of the population, at \texteuro\CapMedian{}, and the Gini coefficient of \CapGiniTracts{} tracts at \CapGini{}. Fourth, outside Navarra, published medians are multiples of \texteuro 350 (\MedBinnedPct{}\% of tracts), consistent with the midpoints of \texteuro 700 intervals, and the P80/P20 ratio is rounded to one decimal.

\paragraph{Other sources} To project incomes beyond \BaseYear{} (Section \ref{sec:nowcast}), I use the annual tax revenue report of the Spanish tax agency (\textit{Informe Anual de Recaudación Tributaria}, Agencia Tributaria). Its Table 2.1 reports household income from tax sources (withholdings and personal income-tax returns, the records the ADRH is built from) by type of income, together with the income tax accrued, about six months after the end of the year; the latest two years are provisional. Population is the national accounts series of Eurostat (\texttt{nama\_10\_pe}). The municipal indicators that \href{https://comparatuingreso.es/}{\mbox{comparatuingreso.es}} shows alongside its results (the share of the population aged 15 and over with higher education and the share born abroad) come from the INE's annual population census and play no role in the estimation.

\paragraph{Territorial units} Spain is divided into 17 autonomous communities, two autonomous cities (Ceuta and Melilla), 52 provinces and, in the ADRH, \MunAdrh{} municipalities. For statistical purposes, municipalities are divided into districts and districts into census tracts (\textit{secciones censales}), which typically have between 1,000 and 2,500 residents.

\paragraph{Summary statistics} Table \ref{tab:summary} describes the \TractsTotal{} tracts in \BaseYear{}. The average tract has \TractPopMean{} residents (standard deviation \TractPopSD{}). Mean net income per consumption unit averages \texteuro\TractEquivMean{} across tracts, \EquivOverPcPct{}\% more than income per person (\texteuro\TractPcMean{}), and the Gini coefficient averages \TractGiniMean{} points. The average tract has a dependency ratio of \TractDepRatioMean{}, a mean age of \TractMeanAge{} years and \TractSingleHhMean{}\% single-person households.

\begin{table}[H]
\centering
\begin{threeparttable}
\onehalfspacing
\caption{Summary statistics, census tracts}\label{tab:summary}
\footnotesize
\begin{tabular}{@{}lccccc@{}}
\toprule
 & Min & Max & Mean & SD & \% missing \\
\midrule
\multicolumn{6}{@{}l}{\textit{Income distribution}} \\
\quad Net income per person (\texteuro) & 5,522 & 36,918 & 15,025 & 4,432 & 1.8 \\
\quad Net income per consumption unit (\texteuro) & 8,797 & 58,752 & 22,036 & 6,915 & 5.5 \\
\quad Median income per consumption unit (\texteuro) & 7,350 & 44,450 & 19,877 & 5,598 & 5.5 \\
\quad Gini index & 19.7 & 44.2 & 28.6 & 3.9 & 5.5 \\
\multicolumn{6}{@{}l}{\textit{Demographics}} \\
\quad Population & 2 & 12,327 & 1,315 & 678 & 1.6 \\
\quad Dependency ratio & 0.00 & 4.00 & 0.63 & 0.18 & 1.6 \\
\quad Mean age & 27.9 & 74.8 & 45.6 & 5.5 & 1.6 \\
\quad Single-person households (\%) & 0.0 & 100.0 & 30.7 & 9.6 & 1.6 \\
\midrule
Census tracts & \multicolumn{5}{l}{37,072} \\
\bottomrule
\end{tabular}

\begin{tablenotes}[para,flushleft]
\scriptsize
\textbf{Notes}: Unweighted statistics across census tracts. Income in \BaseYear{} euros. The dependency ratio is the population under 18 and over 65 relative to the population aged 18 to 65. Source: INE, ADRH, and author's calculations.
\end{tablenotes}
\end{threeparttable}
\end{table}

\paragraph{Coverage} \GiniMissTracts{} tracts (\GiniMissPct{}\%) lack the mean and median per consumption unit, the Gini coefficient and the P80/P20 ratio. They are tracts below the 100-resident threshold: the median one has \GiniMissMedianPop{} residents, whereas every tract with a published Gini coefficient has at least \ObsGiniMinPop{}. Together they house \GiniMissPopPct{}\% of the population. For \GiniImputedTracts{} of them the ADRH publishes income per person. I impute their mean income per consumption unit as income per person times the population-weighted ratio of the two measures in the province, and their Gini coefficient as described in Section \ref{sec:gini}. The remaining \TractsDropped{} tracts have no income data and are dropped. After imputation, the share of tracts without income per consumption unit falls from \MissEquivPct{}\% to \MissEquivAfterPct{}\%.


\section{Methodology}\label{sec:method}


\subsection{Background}

A long-standing empirical regularity is that income distributions are approximately log-normal \citep{aitchison1957lognormal}. Multiplicative processes provide a rationale: incomes result from the interaction of factors such as skills, local labor markets and productivity \citep{neal2000theories}, and if they grow by proportional shocks that are independent of their level, as in Gibrat's law of proportionate effect \citep{gibrat1931inegalites}, the logarithm of income is a sum of independent shocks and tends to normality. The approximation works best within groups exposed to similar shocks. \citet{battistin2009consumption} show that, within cohorts, consumption, which tracks permanent income, is closer to log-normal than current income, whose transitory component departs from log-normality. Mixtures of log-normals fit heterogeneous populations well, each component describing a relatively homogeneous subgroup \citep{flachaire2007, lubrano2016}, and they improve substantially on a single log-normal when mapping income distributions across small areas \citep{gardini2022}.

The literature on neighborhood effects suggests why census tracts are natural components of such a mixture. Contextual influences and endogenous social interactions, such as shared information about jobs, peer effects and common exposure to local shocks, make income-generating processes more alike within neighborhoods than between them \citep{manski1993identification, durlauf1996neighborhoods, durlauf2004neighborhood}.

The log-normal is known to fit the tails poorly. Top incomes follow a power law rather than a log-normal decay \citep{pareto1896cours, atkinson2011top, gabaix2016power}, and the lower tail of Spanish tract distributions is heavier than a log-normal with the same median and Gini coefficient implies. I therefore use a more flexible family, the generalized beta distribution of the second kind (GB2) of \citet{mcdonald1984some}. It nests the log-normal as a limiting case, together with other distributions commonly fitted to incomes, and has a Pareto upper tail \citep{kleiber2003statistical}; and it can be estimated from the kind of summary statistics that statistical offices publish \citep{chotikapanich2007estimating}.

\subsection{Tract distributions and their aggregation}

Let $Y$ be the income per consumption unit of a person living in tract $j$. I assume that $Y$ follows a GB2 distribution, with cumulative distribution function
\begin{equation}\label{eq:gb2}
F_j(y) = I_{t_j(y)}(p_j, q_j), \qquad t_j(y) = \frac{(y/b_j)^{a_j}}{1 + (y/b_j)^{a_j}},
\end{equation}
where $I_t(p, q)$ is the regularized incomplete beta function, $b_j > 0$ is a scale parameter and $a_j, p_j, q_j > 0$ are shape parameters. The GB2 nests the Singh--Maddala ($p_j = 1$) and Dagum ($q_j = 1$) distributions and, as a limiting case, the log-normal. Its upper tail is Pareto with index $a_j q_j$, which I constrain to exceed one so that the mean is finite, and I parameterize it by $\theta_j = (\ln a_j, \ln m_j, \ln p_j, \ln(a_j q_j - 1))$, where $m_j$ is the median.

I estimate $\theta_j$ separately for each tract by weighted minimum distance to the thirteen indicators that the ADRH publishes for it,
\begin{equation}\label{eq:md}
\hat\theta_j = \arg\min_{\theta} \sum_{k=1}^{13} \rho_{jk}\big(g_k(\theta) - g_{jk}\big),
\end{equation}
where $g_{jk}$ is the published value of indicator $k$ and $g_k(\theta)$ its value under the GB2: the median, the mean, the Gini coefficient (computed, as in the ADRH, after replacing the lowest and highest 1.5\% of incomes), the P80/P20 ratio and the nine shares. Differences are measured in log points for the median, the mean and the P80/P20 ratio, in points for the Gini coefficient and in percentage points for the shares, and $\rho_{jk}(d) = (d/s_{jk})^2$, where $s_{jk}$ combines the rounding of the published figure with a tolerance of \TolShare{} points for the shares, \TolGini{} points for the Gini coefficient, \TolMeanPct{}\% for the mean and \TolPPct{}\% for the P80/P20 ratio. Two features of the publication enter $\rho_{jk}$. Medians, published as the midpoints of \texteuro 700 intervals, are treated as intervals: inside the interval the pull towards the published value is weak ($s = \MedPullPct{}\%$) and outside it steep ($s = \MedOutPct{}\%$). Values at the ADRH's caps (Section \ref{sec:data}) are treated as censored: beyond the cap, the tolerance widens to \CapAllowPct{}\% of the capped value. The minimization uses a Levenberg--Marquardt algorithm started from Dagum and Singh--Maddala fits, themselves started from the log-normal with the tract's median and Gini coefficient; \GbConvergedPct{}\% of the \GbTracts{} tracts converge.

The shares are not published for \GbPriorTracts{} of these tracts. For them, the share targets are predictions from regressions of the logit of each published share on a cubic polynomial in the tract's log median, log ratio of mean to median, Gini coefficient and log P80/P20 ratio, and on log population, with province effects shrunk towards zero, estimated on the \GbPriorTrain{} tracts with published shares and no capped indicator. Their tolerance adds the regressions' root mean square error (\GbPriorRMSE{} percentage points on average). In the \GbGapMuns{} municipalities that publish shares, the predictions are shifted so that the municipality's tracts together reproduce its published shares.

The \LocMeanTracts{} tracts without these indicators (Section \ref{sec:data}) get a log-normal, $\ln Y \sim N(\nu_j, \sigma_j^2)$. Its Gini coefficient is $G_j = 2\,\Phi(\sigma_j/\sqrt{2}) - 1$, where $\Phi$ is the standard normal cumulative distribution function, so, given the tract's imputed Gini coefficient $G_j$ (Section \ref{sec:gini}) and imputed mean $\bar y_j$, I set
\begin{equation}\label{eq:params}
\sigma_j = \sqrt{2}\,\Phi^{-1}\!\left(\frac{G_j + 1}{2}\right), \qquad \nu_j = \ln \bar y_j - \frac{\sigma_j^2}{2}.
\end{equation}

The income distribution of an area $A$ (Spain, a province or a municipality) is the population-weighted mixture of the distributions of its tracts, with cumulative distribution function
\begin{equation}\label{eq:mixture}
F_A(y) = \sum_{j \in A} w_j\, F_j(y), \qquad w_j = \frac{N_j}{\sum_{k \in A} N_k},
\end{equation}
where $N_j$ is the population of tract $j$. The $p$-th percentile $q_{A,p}$, $p = 1, \dots, 99$, solves $F_A(q_{A,p}) = p/100$. I solve this equation with Brent's method, bracketing the root on a grid of incomes; the solutions satisfy it to within $10^{\RootErrExp}$. Of the \MunLookup{} municipalities covered, \MunSingleTract{} (\MunSingleTractPct{}\%, home to \MunSingleTractPopPct{}\% of the population) consist of a single tract, so their distribution is that tract's. The \MunDropped{} municipalities in which no tract has income data are not covered.

\subsection{Missing Gini coefficients}\label{sec:gini}

For the \GiniImputedTracts{} tracts with income data but no published Gini coefficient, which get the log-normal, I predict it with gradient-boosted regression trees (XGBoost) trained on the \GiniTrainTracts{} tracts with a published Gini. The predictors are the logarithm of mean income per consumption unit, the dependency ratio, mean age, the share of the population under 18, the share of single-person households, mean household size, population and province dummies. The model uses \XgbRounds{} boosting rounds, a maximum tree depth of \XgbDepth{} and a learning rate of \XgbEta{}, without tuning. Table \ref{tab:gini} compares it with OLS on the same predictors, with province fixed effects, and with a constant, using the same \CvFolds{}-fold cross-validation folds for all three. XGBoost reduces the out-of-sample root mean square error (RMSE) by \XgbRMSEGain{}\% relative to OLS, with an out-of-sample $R^2$ of \XgbRsq{}.

\begin{table}[H]
\centering
\begin{threeparttable}
\onehalfspacing
\caption{Gini imputation: out-of-sample fit}\label{tab:gini}
\footnotesize
\begin{tabular}{@{}lcccc@{}}
\toprule
Model & RMSE & MAE & MAPE (\%) & $R^2$ \\
\midrule
Constant (training mean) & 3.87 & 2.96 & 10.33 & 0.00 \\
OLS, province fixed effects & 2.94 & 2.25 & 7.87 & 0.42 \\
XGBoost, province dummies & 2.57 & 1.97 & 6.89 & 0.56 \\
\bottomrule
\end{tabular}

\begin{tablenotes}[para,flushleft]
\scriptsize
\textbf{Notes}: \CvFolds{}-fold cross-validation over the \GiniTrainTracts{} tracts with a published Gini coefficient, with the same folds for every model. RMSE is the root mean square error, MAE the mean absolute error and MAPE the mean absolute percentage error, all out of sample; $R^2$ is one minus the ratio of the out-of-sample squared error to the variance of the Gini coefficient. The Gini coefficient is measured in points (0--100).
\end{tablenotes}
\end{threeparttable}
\end{table}

Cross-validation measures accuracy for tracts like those in the training sample, all of which have at least \ObsGiniMinPop{} residents, whereas the tracts to impute are smaller. The model extrapolates poorly in that direction: trained on tracts with at least \SmallTrainMin{} residents, it predicts the Gini coefficient of the \SmallTestTracts{} tracts with \ObsGiniMinPop{} to \SmallTestMaxPop{} residents with an RMSE of \XgbSmallRMSE{}, no better than the standard deviation of the Gini across tracts (\GiniSD{}). The imputed Gini coefficients average \GiniPredMean{}, against \GiniObsMean{} for published ones. The imputation hardly affects the national and provincial distributions, since these tracts house \GiniMissPopPct{}\% of the population, but it matters for the smallest municipalities. In \MunSynthetic{} municipalities (\MunSyntheticPct{}\% of those covered, with \MunSyntheticPopPct{}\% of the population), no tract has a published Gini coefficient, so every tract gets the log-normal. Their income per person is, moreover, an average over small municipalities of the same agricultural district, which takes only \MunSyntheticIncomeValues{} distinct values across them. Their distributions are therefore model-based approximations rather than estimates specific to each municipality.

\subsection{Nowcasting}\label{sec:nowcast}

The ADRH is published almost two years after the income year, so its latest edition refers to \BaseYear{}. To compare households with a distribution of the same year, I project incomes to \TargetYear{} with the ECV, whose wave $t+1$ reports mean income per consumption unit in year $t$.

Let $g^{\text{ADRH}}_t$ and $g^{\text{ECV}}_t$ be the growth of national mean income per consumption unit in year $t$ in each source, the ADRH mean being the population-weighted average over the municipalities observed in both years. Survey and administrative growth rates need not coincide, so I measure how ECV growth has compared with ADRH growth over the years both cover,
\begin{equation}
\rho = \frac{\displaystyle\sum_{t=\NowcastFirstYear}^{\NowcastLastYear} g^{\text{ADRH}}_t}{\displaystyle\sum_{t=\NowcastFirstYear}^{\NowcastLastYear} g^{\text{ECV}}_t},
\end{equation}
and nowcast growth in \TargetYear{} as
\begin{equation}
\hat{g}_{\TargetYear} = \rho \, g^{\text{ECV}}_{\TargetYear}.
\end{equation}
Over \NowcastFirstYear--\NowcastLastYear{} the ECV slightly overstates administrative income growth ($\rho = \NowcastRho{}$). With ECV growth of \EcvGrowth{}\% in \TargetYear{}, the nowcast growth is \NowcastGrowth{}\%.

I scale the income of every tract by $1+\hat{g}_{\TargetYear}$. Multiplying a GB2 variable by a constant multiplies its scale $b_j$ by that constant and leaves $a_j$, $p_j$ and $q_j$ unchanged (for a log-normal, it shifts $\nu_j$ by $\ln(1+\hat{g}_{\TargetYear})$ and leaves $\sigma_j$ unchanged), so the adjustment moves every distribution while preserving its shape and Gini coefficient. The underlying assumption is that incomes grew at a common rate across areas and along the distribution over one year.

The projection is uncertain. In leave-one-year-out validation, $\rho\, g^{\text{ECV}}_t$ misses $g^{\text{ADRH}}_t$ by \NowcastLooMAE{} percentage points on average, against \EcvRawMAE{} points for uncorrected ECV growth. The estimate of $\rho$ also depends on the window and the estimator: it is \RhoLastThree{} over the last three years, \RhoExclCovid{} excluding 2020 and 2021, and \RhoOrigin{} by least squares through the origin. Across these choices, including no correction, nowcast growth ranges from \NowcastGrowthMin{}\% to \NowcastGrowthMax{}\%.

\subsection{Placing a household}\label{sec:placing}

\href{https://comparatuingreso.es/}{\mbox{Comparatuingreso.es}} asks for the household's net monthly income in \TargetYear{}, the year of the nowcast distribution, and for whether it was paid in 12 or 14 instalments (Spanish salaries and pensions are often paid in 14, with two extra payments a year). Annual income is the monthly amount times 12 or 14. Dividing it by the household's consumption units on the modified OECD scale, from the number of members aged 14 or over and under 14, gives income per consumption unit, $y$.

The household's position in area $A$ is
\begin{equation}
P_A(y) =
\begin{cases}
0 & \text{if } y < q_{A,1}, \\
\max\{p \in \{1, \dots, 99\} : q_{A,p} \le y\} & \text{otherwise.}
\end{cases}
\end{equation}
Since $F_A(y) \ge P_A(y)/100$, a household with $P_A(y) \ge 1$ has more income per consumption unit than at least $P_A(y)\%$ of the population, which is how the app reports it; $P_A(y) = 0$ is reported as being among the 1\% with the lowest income. Positions are shares of persons, not households: every person counts with the income per consumption unit of their household.


\section{Results}\label{sec:results}


\subsection{Fit at the tract level}

The tract distributions are fitted to every indicator the ADRH publishes for the tract, so comparing them with those indicators measures the quality of the fit rather than its predictive accuracy. For the \ValTracts{} tracts with a GB2 fit, all averages and regressions below weight tracts by population. The fitted median lies inside the published \texteuro 700 interval in \FitMedInBinPct{}\% of tracts; fitted and published means differ by \FitMeanAbsErr{}\% on average, Gini coefficients by \FitGiniAbsErr{} points and P80/P20 ratios by \FitPAbsErr{}\%, the latter partly because published ratios come from rounded percentiles. The nine shares differ from the published ones by \FitShareMAE{} percentage points on average.

\paragraph{Mean-to-median ratio} The ratio of the mean to the median summarizes the skewness of each tract's distribution. Figure \ref{fig:skew} plots the observed $\ln(\bar y_j / \tilde m_j)$, where $\bar y_j$ and $\tilde m_j$ are the published mean and median, against its fitted value. A regression of observed on fitted values has a slope of \ValSkewSlope{} ($R^2 = \ValSkewRsq{}$).

\begin{figure}[!htb]
\begin{center}
\caption{Mean-to-median ratio: observed and fitted}
\label{fig:skew}
\includegraphics[width=0.55\textwidth]{output/fig_validation_skew.png}
\end{center}
\begin{fignotes2}
\textbf{Notes:} Binned scatter plot of the observed $\ln(\text{mean}/\text{median})$ of income per consumption unit against its value under the tract's fitted GB2. Tracts are grouped into 30 bins of equal population. The solid line is the population-weighted regression of observed on fitted values over all tracts; the dashed line is the 45-degree line. Source: INE, ADRH (\BaseYear{}), and author's calculations.
\end{fignotes2}
\end{figure}

\paragraph{P80/P20 ratio} Figure \ref{fig:p80p20} does the same for the P80/P20 ratio. A regression of observed on fitted ratios has a slope of \ValPSlope{} ($R^2 = \ValPRsq{}$).

\begin{figure}[!htb]
\begin{center}
\caption{P80/P20 ratio: observed and fitted}
\label{fig:p80p20}
\includegraphics[width=0.55\textwidth]{output/fig_validation_p80p20.png}
\end{center}
\begin{fignotes2}
\textbf{Notes:} Binned scatter plot of the observed P80/P20 ratio of income per consumption unit against its value under the tract's fitted GB2. Tracts are grouped into 30 bins of equal population. The solid line is the population-weighted regression of observed on fitted values over all tracts; the dashed line is the 45-degree line. Source: INE, ADRH (\BaseYear{}), and author's calculations.
\end{fignotes2}
\end{figure}

\paragraph{Held-out shares} To measure how well the GB2 predicts what it is not fitted to, I refit every tract nine times, each time leaving one share out of its targets, and compare the predicted share with the published one over the tracts that publish it. The held-out shares differ from the published ones by \HoldShareMAE{} percentage points on average, between \HoldShareMAEMin{} and \HoldShareMAEMax{} depending on the share (Table \ref{tab:thresholds}, last column), against \FitShareMAE{} points when the share is part of the fit.

\subsection{Validation of the aggregate distributions}

What the app reports are percentiles of the mixtures. The mean of a mixture and its shares below income thresholds are population-weighted averages of those of its tracts, so they largely reproduce the tract fits. Its median, Gini coefficient and P80/P20 ratio are not, and the ADRH publishes them for Spain, its provinces and its municipalities; they are the main out-of-sample test of the model. Table \ref{tab:aggregate} compares them with the published values; the comparison is in \BaseYear{} euros, before the nowcast, which scales every income by the same factor.

The national median of the mixture, \texteuro\MixNatMedian{}, lies inside the published interval (\texteuro\AdrhNatMedianLo{} to \texteuro\AdrhNatMedianHi{}), \MixNatMedianErr{}\% from its midpoint. The mean, \texteuro\MixNatMean{}, differs by \MixNatMeanErr{}\% from the published \texteuro\AdrhNatMean{}. Computed as the ADRH computes it, after replacing the lowest and highest 1.5\% of incomes, the Gini coefficient of the mixture is \MixNatGini{}, against \AdrhNatGini{} published. Its P80/P20 ratio is \MixNatPEightyTwenty{}, or \MixNatPEightyTwentyBinned{} once its 20th and 80th percentiles are rounded to \texteuro 700 intervals as the published ones are, against \AdrhNatPEightyTwenty{}. The model median lies inside the published interval for \ProvMedInBinPct{}\% of provinces and for \MunMedInBinPct{}\% of the \MunValN{} municipalities with at least two tracts and published statistics (population-weighted). Gini coefficients differ from the published ones by \ProvGiniAbsErr{} points in provinces and \MunGiniAbsErr{} points in municipalities on average, and P80/P20 ratios by \ProvPAbsErr{}\% and \MunPAbsErr{}\%.

\begin{table}[H]
\centering
\begin{threeparttable}
\onehalfspacing
\caption{Model and published statistics}\label{tab:aggregate}
\footnotesize
\begin{tabular}{@{}lccc@{}}
\toprule
\textit{A. Spain} & ADRH & Model & Difference (\%) \\
\midrule
\quad Median (\texteuro) & 19,250 & 19,029 & $-$1.1 \\
\quad Mean (\texteuro) & 22,318 & 22,326 & $+$0.04 \\
\quad Gini index & 32.1 & 32.2 & $+$0.2 \\
\quad P80/P20 & 2.8 & 2.71 & $-$3.4 \\
\midrule
\textit{B. Provinces and municipalities} & Units & Mean abs.\ diff. & Mean diff. \\
\midrule
\quad Provinces: median (\%) & 52 & 0.8 & $+$0.2 \\
\quad Provinces: Gini index (points) & 52 & 0.07 & $+$0.03 \\
\quad Provinces: P80/P20 (\%) & 52 & 1.5 & $-$0.3 \\
\quad Municipalities: median (\%) & 2,253 & 0.9 & $+$0.1 \\
\quad Municipalities: Gini index (points) & 2,253 & 0.11 & $+$0.02 \\
\quad Municipalities: P80/P20 (\%) & 2,253 & 2.0 & $+$0.5 \\
\bottomrule
\end{tabular}

\begin{tablenotes}[para,flushleft]
\scriptsize
\textbf{Notes}: Statistics of income per consumption unit, in \BaseYear{} euros. Panel A compares the mixture for Spain with the totals published by the ADRH; the Gini coefficient of the mixture is computed after replacing the lowest and highest 1.5\% of incomes with the nearest remaining value, as in the ADRH. Published medians are midpoints of \texteuro 700 intervals and published P80/P20 ratios are rounded to one decimal. Panel B reports population-weighted averages of the difference between model and published statistics, in absolute value and with its sign, in per cent for medians and P80/P20 ratios and in points for Gini coefficients, for the provinces and for the municipalities with at least two census tracts and published statistics. Source: INE, ADRH, and author's calculations.
\end{tablenotes}
\end{threeparttable}
\end{table}

Table \ref{tab:thresholds} compares the shares of the population below nine thresholds that the ADRH publishes for Spain (three fixed amounts and six fractions of the national median) with those of the mixture. They differ by at most \ThrOtherMaxDiff{} percentage points, except below \texteuro 5,000, where the mixture places \ThrFiveKModel{}\% of the population against \ThrFiveKObs{}\% in the ADRH. The share of the population with the lowest incomes is thus somewhat understated: the tract GB2s cannot fully reproduce the mass of very low incomes while also matching the shares higher up.

\begin{table}[H]
\centering
\begin{threeparttable}
\onehalfspacing
\caption{Share of the population below income thresholds}\label{tab:thresholds}
\footnotesize
\begin{tabular}{@{}lcccccc@{}}
\toprule
 & & \multicolumn{3}{c}{Spain: share below (\%)} & \multicolumn{2}{c}{Tracts: MAE (pp)} \\ \cmidrule(lr){3-5} \cmidrule(l){6-7}
Threshold & \texteuro{} (2023) & ADRH & Model & Difference (pp) & Fitted & Held out \\
\midrule
Fixed amount & 5,000 & 4.3 & 3.9 & $-$0.4 & 0.69 & 0.87 \\
Fixed amount & 7,500 & 9.0 & 9.0 & 0.0 & 0.45 & 0.55 \\
40\% of the median & 7,700 & 9.5 & 9.5 & 0.0 & 0.48 & 0.58 \\
50\% of the median & 9,625 & 14.7 & 14.9 & $+$0.2 & 0.46 & 0.55 \\
Fixed amount & 10,000 & 15.8 & 16.1 & $+$0.3 & 0.47 & 0.58 \\
60\% of the median & 11,550 & 21.3 & 21.3 & 0.0 & 0.59 & 0.84 \\
140\% of the median & 26,950 & 73.6 & 73.5 & $-$0.1 & 0.46 & 0.89 \\
160\% of the median & 30,800 & 80.9 & 80.9 & 0.0 & 0.50 & 0.76 \\
200\% of the median & 38,500 & 90.0 & 89.8 & $-$0.2 & 0.47 & 0.73 \\
\bottomrule
\end{tabular}

\begin{tablenotes}[para,flushleft]
\scriptsize
\textbf{Notes}: Shares of the population with income per consumption unit below each threshold. The relative thresholds are fractions of the national median (\texteuro\AdrhNatMedian{}); for those above the median the ADRH publishes the share above, shown here as its complement. The first columns compare the shares published for Spain with those of the mixture. The last two columns are population-weighted mean absolute differences between the shares of the tracts' fitted GB2s and those the ADRH publishes for the tracts, over the tracts that publish them (\ShareTractsPopPct{}\% of the population): with every share in the fit (fitted) and with the share left out of the fit (held out). Source: INE, ADRH, and author's calculations.
\end{tablenotes}
\end{threeparttable}
\end{table}

For households, the first percentile of the mixture is \texteuro\MixPOne{}, and a household at the at-risk-of-poverty line (60\% of the median, \texteuro\ThrPovEuro{}) is placed at position \ThrPovModelP{} (Section \ref{sec:placing}), against \ThrPovObsP{} in the published distribution. At twice the median (\texteuro\ThrTopEuro{}), the mixture places \ThrTopModel{}\% of the population above, against \ThrTopObs{}\% in the ADRH.

\subsection{Decomposing income inequality}

To examine the spatial structure of inequality, I decompose the variance of log income within each autonomous community $c$ into a within-tract component and between-area components:
\begin{equation}\label{decomp}
\sigma^2_{c} =
\underbrace{\sum_{j \in c} w_j \sigma^2_{j}}_{\text{Within tracts}} +
\underbrace{\sum_{j \in c} w_j \big( \nu_{j} - \bar\nu_{k(j)} \big)^2}_{\text{Between tracts}} +
\underbrace{\sum_{k \in c} w_k \big( \bar\nu_{k} - \bar\nu_{p(k)} \big)^2}_{\text{Between municipalities}} +
\underbrace{\sum_{p \in c} w_p \big( \bar\nu_{p} - \bar\nu_{c} \big)^2}_{\text{Between provinces}},
\end{equation}
where $\nu_j$ and $\sigma_j^2$ are the mean and variance of log income in tract $j$, which for a GB2 are $\nu_j = \ln b_j + [\psi(p_j) - \psi(q_j)]/a_j$ and $\sigma_j^2 = [\psi'(p_j) + \psi'(q_j)]/a_j^2$, with $\psi$ the digamma function; $\bar\nu_k$, $\bar\nu_p$ and $\bar\nu_c$ are the population-weighted means of $\nu_j$ in municipality $k$, province $p$ and community $c$, $k(j)$ is the municipality of tract $j$ and $p(k)$ the province of municipality $k$, and $w_j$, $w_k$ and $w_p$ are population shares within $c$.

Figure \ref{fig:variance_decomp} shows the result for each community. Within-tract variation accounts for between \VdCcaaWithinMin{}\% and \VdCcaaWithinMax{}\% of the variance of log income, and for \VdCcaaWithin{}\% on average across the \VdCcaaN{} communities.\footnote{This does not contradict the homogeneity of income-generating processes within neighborhoods discussed in Section \ref{sec:method}, which refers to the mechanisms that shape incomes, such as proportional shocks, and not to the level of incomes.} Differences between tracts of the same municipality account for a further \VdCcaaTract{}\% on average, and differences between municipalities and provinces for \VdCcaaRest{}\%. For Spain as a whole, within-tract variation accounts for \VdWithin{}\% of the variance, differences between tracts within municipalities for \VdBetweenTract{}\%, between municipalities within provinces for \VdBetweenMun{}\%, between provinces within communities for \VdBetweenProv{}\% and between communities for \VdBetweenCcaa{}\%. Most inequality therefore lies within tracts, which is why the aggregate distribution depends on how incomes are distributed inside each tract and not only on tract averages. These shares rest on the fitted GB2s.

\begin{figure}[H]
\begin{center}
\caption{Decomposition of the variance of log income by autonomous community}
\label{fig:variance_decomp}
\includegraphics[width=0.85\textwidth]{output/fig_variance_decomp.png}
\end{center}
\begin{fignotes2}
\textbf{Notes:} Decomposition of equation (\ref{decomp}) for each autonomous community, excluding Ceuta and Melilla. Tracts with a GB2 fit. Source: INE, ADRH (\BaseYear{}), and author's calculations.
\end{fignotes2}
\end{figure}


\section{Conclusion}\label{sec:conc}

This note has described a transparent way to estimate income distributions for Spain, its provinces and its municipalities from tract-level statistics: fit a GB2 to the thirteen statistics published for each tract, and aggregate tracts into population-weighted mixtures. Out of sample, the resulting distributions reproduce published medians, Gini coefficients and P80/P20 ratios closely at every level, and the national shares below income thresholds within \ThrOtherMaxDiff{} percentage points above \texteuro 5,000.

Their main limitation is the lowest incomes: the model places \ThrFiveKModel{}\% of the population below \texteuro 5,000 per consumption unit, against \ThrFiveKObs{}\%, and so slightly understates the position of the poorest households. Other limitations are smaller or affect few users. The top of the distribution rests on the GB2's Pareto tail beyond the highest published threshold and on tract medians that the ADRH caps. The distributions of the smallest municipalities rest on imputed inputs and a log-normal. The two-year nowcast has missed the published income level by \NowcastBtMAE{} percentage points on average in the past. Finally, the ADRH measures income recorded by the tax authorities, which includes remuneration in kind and notional income from second homes and excludes informal earnings and transfers between households, whereas users report the money that reaches their household.

Natural extensions are a component for very low incomes within each tract, to close the remaining gap below \texteuro 5,000, and the use of external information on top incomes to discipline the upper tail. The framework can be adapted to other countries with comparable small-area data.

%Bibliography
\newpage
\singlespacing
\bibliography{bib.bib}

\end{document}
